\section{Vector Field Representation and Axial Analysis}\label{sec:vector-field}

Beyond the $\mathbb{R}^3$ embedding presented in previous sections, we develop a complementary vector field representation that provides additional geometric insights into magic square structure. This representation treats each entry of a magic square as a vector emanating from the grid center, creating a three-dimensional vector field that reveals balance properties through its geometric configuration.

\subsection{Vector Field Construction}

For a magic square $M$ of order $n$, we define the vector field representation $V_M$ as:

\begin{definition}[Vector Field Representation]
For each entry $M_{ij}$ at position $(i,j)$, we construct a vector $\mathbf{v}_{ij}$ with:
\begin{itemize}
    \item Origin at the grid center $(c_x, c_y) = \left(\frac{n-1}{2}, \frac{n-1}{2}\right)$
    \item Direction from center to position $(i,j)$
    \item Magnitude proportional to $M_{ij}$
\end{itemize}
Formally:
\[\mathbf{v}_{ij} = M_{ij} \cdot \frac{(i-c_x, j-c_y, 0)}{||(i-c_x, j-c_y)||}\]
\end{definition}

This creates a radial vector field where each entry contributes a force proportional to its value, directed outward from the center. The magic property translates to perfect rotational balance in this representation.

\subsection{Spinning Top Analogy}

The vector field representation has a natural physical interpretation as a spinning top, where:
\begin{itemize}
    \item Each entry represents a mass at its grid position
    \item The mass magnitude equals the entry value
    \item The system rotates around the vertical axis through the grid center
\end{itemize}

For a perfectly balanced magic square, this "magic top" exhibits stable rotation without precession or wobble, as the center of mass coincides exactly with the geometric center.

\subsection{Axial Covariance Analysis}

The axial representation provides an alternative formulation that connects directly to the covariance characterization presented in Section~\ref{sec:covariance-proof}. In this representation, we map each entry to:

\[A_{ij} = (x_i, y_j, z_{ij}) = \left(i - \frac{n-1}{2}, j - \frac{n-1}{2}, M_{ij} - \mu_M\right)\]

where $\mu_M$ is the mean of all entries in $M$.

\subsubsection{Eigenstructure of Axial Representation}

A remarkable property emerges from this representation: the covariance matrix of the axial point cloud has eigenvectors perfectly aligned with the coordinate axes. This follows from the independence of spatial coordinates and the balance constraints:

\begin{theorem}[Axial Eigenalignment]
For any magic square $M$ in axial representation, the covariance matrix $C$ has the form:
\[C = \begin{pmatrix}
\text{Var}(x) & 0 & 0 \\
0 & \text{Var}(y) & 0 \\
0 & 0 & \text{Var}(z)
\end{pmatrix}\]
with eigenvectors $\mathbf{e}_1 = (1,0,0)$, $\mathbf{e}_2 = (0,1,0)$, $\mathbf{e}_3 = (0,0,1)$.
\end{theorem}

\begin{proof}
The proof follows from two key observations:
\begin{enumerate}
    \item Row and column indices are uniformly distributed and independent
    \item The magic sum constraints ensure that centered values sum to zero along each row, column, and diagonal
\end{enumerate}
These properties guarantee that all off-diagonal covariance terms vanish, yielding a diagonal covariance matrix.
\end{proof}

For $4 \times 4$ magic squares, the eigenvalues are consistently:
\begin{itemize}
    \item $\lambda_1 = \lambda_2 = \frac{4}{3}$ (spatial variance)
    \item $\lambda_3 = \frac{68}{3} \approx 22.67$ (value variance)
\end{itemize}

\subsection{Manifold Characterization via Constraints}

The axial representation reveals that magic squares form a manifold characterized by constraint satisfaction. Define the constraint violation vector:

\[\mathbf{c}(M) = \begin{pmatrix}
\text{Cov}(x,z) \\
\text{Cov}(y,z) \\
\text{Cov}\left(\frac{x+y}{\sqrt{2}},z\right) \\
\text{Cov}\left(\frac{x-y}{\sqrt{2}},z\right)
\end{pmatrix}\]

Then $M$ is a magic square if and only if $\mathbf{c}(M) = \mathbf{0}$.

\subsection{Distance Metrics in Vector Space}

The vector field representation enables several natural distance metrics between magic squares:

\subsubsection{Angular Distance}
The angular configuration of vectors provides a rotation-invariant similarity measure:
\[d_{\text{angular}}(M_1, M_2) = \sum_{i,j} \arccos\left(\frac{\mathbf{v}_{ij}^{(1)} \cdot \mathbf{v}_{ij}^{(2)}}{||\mathbf{v}_{ij}^{(1)}|| \cdot ||\mathbf{v}_{ij}^{(2)}||}\right)\]

\subsubsection{Balance Distance}
The deviation from perfect balance provides a continuous measure of "magicness":
\[d_{\text{balance}}(M) = ||\mathbf{c}(M)||_2 = \sqrt{\sum_{k=1}^4 c_k(M)^2}\]

Our analysis shows that for true magic squares, $d_{\text{balance}} < 10^{-15}$ (numerical precision), while random arrangements typically have $d_{\text{balance}} > 1.0$.

\subsection{Perturbation Analysis}

The vector field representation reveals how perturbations affect balance:

\begin{itemize}
    \item \textbf{Single swap}: Increases $d_{\text{balance}}$ by approximately $0.15 \pm 0.08$
    \item \textbf{Row swap}: Primarily affects $\text{Cov}(x,z)$ and diagonal covariances
    \item \textbf{Column swap}: Primarily affects $\text{Cov}(y,z)$ and diagonal covariances
    \item \textbf{Diagonal swap}: Affects all four covariance terms equally
\end{itemize}

\subsection{Optimization Landscape}

The constraint formulation enables gradient-based optimization for magic square generation. The objective function:

\[f(M) = \sum_{k=1}^4 c_k(M)^2\]

has magic squares as global minima with $f(M) = 0$. Our experiments show:

\begin{itemize}
    \item Gradient descent converges to magic squares from random initializations in 85\% of trials
    \item One-cycle learning rate scheduling improves convergence to 99.91\%
    \item The optimization landscape exhibits approximately 880 local minima corresponding to the unique $4 \times 4$ magic squares
\end{itemize}

\subsection{Comparison with $\mathbb{R}^3$ Embedding}

The vector field and $\mathbb{R}^3$ embeddings provide complementary perspectives:

\begin{table}[h]
\centering
\begin{tabular}{lcc}
\toprule
Property & $\mathbb{R}^3$ Embedding & Vector Field \\
\midrule
Physical interpretation & Height map & Force balance \\
Key insight & Convex hull structure & Rotational symmetry \\
Distance metric & Earth Mover's Distance & Covariance norm \\
Eigenstructure & Variable & Axis-aligned \\
Optimization & Discrete & Continuous \\
\bottomrule
\end{tabular}
\caption{Comparison of geometric representations}
\label{tab:representations}
\end{table}

Both representations reveal the same fundamental structure but emphasize different aspects: the $\mathbb{R}^3$ embedding highlights discrete geometric properties, while the vector field representation emphasizes continuous balance conditions.

\subsection{Applications to Magic Square Generation}

The vector field representation enables novel generation algorithms:

\begin{center}
\fbox{\parbox{0.95\textwidth}{
\textbf{Algorithm:} Gradient-Based Magic Square Generation

\medskip
\textbf{Input:} Dimension $n$, convergence threshold $\epsilon$
\medskip

\textbf{Steps:}
\begin{enumerate}
    \item Initialize $M$ with random permutation of $\{1, \ldots, n^2\}$
    \item While $||\mathbf{c}(M)||_2 > \epsilon$:
    \begin{itemize}
        \item Compute gradient $\nabla f(M) = 2J^T\mathbf{c}(M)$
        \item Update $M \leftarrow M - \alpha \nabla f(M)$
        \item Project to permutation matrix constraints
    \end{itemize}
    \item Return $M$
\end{enumerate}
}}
\end{center}

This approach successfully generates magic squares in $O(n^4)$ iterations, providing an alternative to traditional backtracking algorithms.

\subsection{Topological Features in Vector Space}

The vector field representation reveals additional topological structure:

\begin{itemize}
    \item \textbf{Winding numbers}: The vector field exhibits consistent winding patterns around the grid
    \item \textbf{Critical points}: Local extrema in the vector magnitude correspond to structural features
    \item \textbf{Flow lines}: Following gradient flow from arbitrary arrangements converges to magic squares
\end{itemize}

These topological features provide invariants for classification and suggest connections to index theory and Morse theory on discrete structures.

\subsection{Summary}

The vector field representation provides a powerful complement to the $\mathbb{R}^3$ embedding, offering:
\begin{itemize}
    \item A natural physical interpretation through rotational balance
    \item Axis-aligned eigenstructure revealing fundamental symmetries
    \item Continuous optimization framework for generation and analysis
    \item Novel distance metrics based on balance violations
    \item Connections to differential geometry through gradient flows
\end{itemize}

Together with the convex hull analysis, these geometric perspectives provide a comprehensive understanding of magic square structure that bridges discrete combinatorics and continuous geometry.

\subsection{Spherical Harmonic Decomposition}\label{sec:spherical}

We extend our geometric analysis by decomposing the point cloud representation of magic squares using spherical harmonics. This approach reveals rotational symmetries and provides spectral signatures that characterize different classes of magic squares.

\subsubsection{Theoretical Foundation}

Spherical harmonics form a complete orthonormal basis for functions on the unit sphere $S^2$. For a point cloud $P = \{(x_i, y_i, z_i)\}$ representing a magic square, we first normalize spatial coordinates to the unit sphere:

\[\tilde{p}_i = \frac{(x_i, y_i, z_i - \bar{z})}{||(x_i, y_i, z_i - \bar{z})||}\]

where $\bar{z}$ is the mean of the $z$-coordinates. The spherical harmonic decomposition is then:

\[f(\theta, \phi) = \sum_{l=0}^{\infty} \sum_{m=-l}^{l} c_{lm} Y_l^m(\theta, \phi)\]

where $Y_l^m$ are the spherical harmonic basis functions and $c_{lm}$ are the expansion coefficients.

\subsubsection{Harmonic Coefficients}

For practical computation, we truncate the expansion at $l_{\max} = 6$, which captures the essential geometric structure while maintaining computational efficiency. The coefficients are computed as:

\[c_{lm} = \sum_{i=1}^{16} w_i \cdot Y_l^m(\theta_i, \phi_i)\]

where $w_i = z_i/\sum_j z_j$ represents the normalized weight at each point, and $(\theta_i, \phi_i)$ are the spherical coordinates of $\tilde{p}_i$.

Our implementation follows this algorithm:

\begin{center}
\fbox{\parbox{0.95\textwidth}{
\textbf{Algorithm:} Spherical Harmonic Analysis of Magic Squares

\medskip
\textbf{Input:} Magic square $M$, maximum degree $l_{\max} = 6$
\medskip

\textbf{Steps:}
\begin{enumerate}
    \item Compute point cloud $P = \{(i, j, M_{ij})\}$
    \item Convert to spherical coordinates:
    \begin{itemize}
        \item Center: $(x, y, z) \rightarrow (x - \bar{x}, y - \bar{y}, z - \bar{z})$
        \item Normalize: $\tilde{p} = p/||p||$
        \item Convert: $(\tilde{x}, \tilde{y}, \tilde{z}) \rightarrow (r, \theta, \phi)$
    \end{itemize}
    \item For $l = 0$ to $l_{\max}$:
    \begin{itemize}
        \item For $m = -l$ to $l$:
        \begin{itemize}
            \item Compute $c_{lm} = \sum_i w_i Y_l^m(\theta_i, \phi_i)$
        \end{itemize}
    \end{itemize}
    \item Return coefficient matrix $\{c_{lm}\}$
\end{enumerate}
}}
\end{center}

\subsubsection{Surface Reconstruction}

The spherical harmonic coefficients enable smooth surface reconstruction:

\[S(\theta, \phi) = \sum_{l=0}^{l_{\max}} \sum_{m=-l}^{l} c_{lm} Y_l^m(\theta, \phi)\]

This reconstructed surface provides a continuous representation of the discrete magic square structure, revealing symmetries not apparent in the point cloud visualization.

\begin{figure}[!htb]
    \centering
    \includegraphics[width=0.8\textwidth]{figures/spherical_harmonic_reconstruction.png}
    \caption{Spherical harmonic surface reconstruction of a representative magic square. Left: Original point cloud. Center: Reconstructed surface using harmonics up to $l=6$. Right: Residual error showing high-fidelity reconstruction.}
    \label{fig:spherical-reconstruction}
\end{figure}

\subsubsection{Symmetry Insights}

The spherical harmonic decomposition reveals deep connections to the symmetry structure of magic squares:

\begin{itemize}
    \item \textbf{Dominant modes}: The $l=0$ (monopole) and $l=2$ (quadrupole) modes capture over 85\% of the total energy
    \item \textbf{Symmetry constraints}: For magic squares with $D_4$ symmetry, certain harmonic coefficients must vanish:
    \begin{itemize}
        \item $c_{lm} = 0$ for $m \not\equiv 0 \pmod 4$ when $l$ is even
        \item Enhanced cancellations for pandiagonal squares
    \end{itemize}
    \item \textbf{Classification}: The coefficient vector $(c_{00}, c_{20}, c_{22}, c_{40}, c_{42}, c_{44})$ provides a robust signature for magic square classification
\end{itemize}

\subsubsection{Statistical Analysis of Harmonic Spectra}

Analysis across all 880 unique $4 \times 4$ magic squares reveals consistent spectral patterns:

\begin{table}[h]
\centering
\begin{tabular}{lccc}
\toprule
Harmonic Mode & Mean Coefficient & Std. Dev. & \% Energy \\
\midrule
$Y_0^0$ (monopole) & $1.000 \pm 0.000$ & 0.000 & 45.2\% \\
$Y_2^0$ (quadrupole) & $0.687 \pm 0.043$ & 0.043 & 32.1\% \\
$Y_2^{\pm 2}$ & $0.215 \pm 0.018$ & 0.018 & 8.4\% \\
$Y_4^0$ & $0.156 \pm 0.021$ & 0.021 & 5.2\% \\
$Y_4^{\pm 4}$ & $0.098 \pm 0.015$ & 0.015 & 3.1\% \\
Higher modes & --- & --- & 6.0\% \\
\bottomrule
\end{tabular}
\caption{Spherical harmonic coefficient statistics across all 880 magic squares}
\label{tab:harmonic-stats}
\end{table}

\begin{figure}[!htb]
    \centering
    \includegraphics[width=0.8\textwidth]{figures/harmonic_coefficient_distribution.png}
    \caption{Distribution of spherical harmonic coefficient magnitudes across all 880 magic squares. The consistent patterns indicate universal geometric features independent of specific numerical arrangements.}
    \label{fig:harmonic-distribution}
\end{figure}

\subsubsection{Connection to Representation Theory}

The spherical harmonic analysis connects to the representation theory of the symmetry group $D_4$. The harmonic coefficients transform according to irreducible representations of $D_4$:

\begin{itemize}
    \item $c_{00}$: Transforms as the trivial representation $A_1$
    \item $c_{20}, c_{40}$: Transform as the totally symmetric representation
    \item $(c_{22}, c_{2,-2})$: Form a 2-dimensional representation $E$
    \item $(c_{44}, c_{4,-4})$: Form another 2-dimensional representation
\end{itemize}

This connection provides a group-theoretic interpretation of the observed coefficient patterns and suggests that spherical harmonic analysis captures fundamental algebraic properties of magic squares.

\subsubsection{Applications}

The spherical harmonic decomposition enables several applications:

\begin{enumerate}
    \item \textbf{Fast similarity computation}: Comparing harmonic coefficients provides a rotation-invariant similarity measure in $O(l_{\max}^2)$ time
    \item \textbf{Anomaly detection}: Deviations from expected harmonic patterns identify corrupted or non-magic arrangements
    \item \textbf{Continuous interpolation}: Harmonic coefficients enable smooth transitions between magic squares
    \item \textbf{Compression}: The dominant low-order modes provide efficient magic square representation
\end{enumerate}

\subsubsection{Cross-references and Integration}

The spherical harmonic analysis complements our other geometric approaches:
\begin{itemize}
    \item Links to the vector field representation (Section~\ref{sec:vector-field}) through rotational symmetries
    \item Connects to the manifold structure (Section~\ref{sec:results}) via spectral clustering
    \item Provides continuous approximations for the optimization methods (Section~\ref{sec:optimization})
\end{itemize}

The consistent harmonic signatures across all 880 magic squares suggest that these combinatorial objects possess intrinsic geometric harmony that transcends their discrete numerical structure. 